\begin{definition}\label{def:IsParadoxical}\lean{IsParadoxical}\leanok
A pair $(j, n)$ consisting of a positive iteration count $j$ and a starting integer $n$ is called \emph{paradoxical} if:
\begin{equation}
    T^j(n) \ge n \quad \text{and} \quad C_j(n) < 1.
\end{equation}
That is, the $j$-th iterate of the Terras map is at least as large as the starting value, and the
decomposition correction at $j$ is smaller than~$1$.
The sequence of values
\[
n, T(n), T^2(n), \cdots, T^{j}(n)
\]
is called a paradoxical sequence.
\end{definition}

\begin{lemma}\label{lemma:E_shift_mul}\lean{E_shift_mul}\leanok
For any $k, j, n \in \mathbb{N}$, the correction ratio behaves under powers of two as:
\begin{equation}
    E_{k+j}(2^k n) = E_j(n).
\end{equation}
\end{lemma}
\begin{proof}\leanok
This follows from the shift properties of the decomposition formula: the number of odd steps $q(k+j, 2^k n) = q(j, n)$ and the decomposition correction $\mathcal{Q}_{k+j}(2^k n) = 2^k \mathcal{Q}_j(n)$. Substituting these into the definition of $E$ clears the $2^k$ factor.
\end{proof}

\begin{lemma}\label{lemma:num_odd_steps_unbounded}\lean{num_odd_steps_unbounded}\leanok
If $n \ge 1$ has an infinite stopping time, then the number of odd steps $q(j, n)$ is unbounded as $j \to \infty$.
\end{lemma}
\begin{proof}\leanok
Suppose $q(j, n)$ were bounded by some $M$. Since $q$ is non-decreasing, it must eventually stay constant at some value $q_s$. This would imply that all iterates $T^j(n)$ for $j \ge J$ are even, so they are repeatedly divided by 2. This contradicts the infinite stopping time, as the sequence would eventually drop below $n$ (or reach 1 and then cycle).
\end{proof}

\begin{lemma}\label{lemma:num_odd_steps_hits_all}\lean{num_odd_steps_hits_all}\leanok
If $n \ge 1$ has an infinite stopping time, then for every $a \in \mathbb{N}$, there exists an iteration count $j$ such that $q(j, n) = a$.
\end{lemma}
\begin{proof}\leanok
Since $q(0, n) = 0$, $q(j, n)$ is unbounded, and $q(j+1, n) - q(j, n) \in \{0, 1\}$, the discrete intermediate value theorem implies that $q$ must hit every natural number.
\end{proof}

\begin{theorem}[Rozier--Terracol Theorem 3.2]\label{theorem:CRoz_lemma_32}\lean{CRoz_lemma_32}\leanok
If the integer $n$ has an infinite stopping time, then there exist infinitely many paradoxical sequences starting from integers of the form $2^k n$ with $k \ge 0$.
\end{theorem}
\begin{proof}\leanok
We construct an infinite set of paradoxical pairs $(j', n')$ where $n'$ is of the form $2^k n$.
Consider the set of Diophantine approximations $S = \{ (a, b) \in \mathbb{N}^2 \mid 1 - \frac{1}{4n} < \frac{3^a}{2^b} < 1 \}$. By Theorem~\ref{theorem:exists_infinite_pairs_approx}, this set is infinite.
For each $(a, b) \in S$, we use Lemma~\ref{lemma:num_odd_steps_hits_all} to find an iteration $j$ such that $q(j, n) = a$. Let $j_{of}(a)$ be the smallest such $j$. We distinguish two cases for each $(a, b) \in S$:

\textbf{Case 1: $j_{of}(a) \ge b$.}
We claim that $(j_{of}(a), n)$ is paradoxical.
First, $T^{j_{of}(a)}(n) \ge n$ because $n$ has infinite stopping time.
Second, $C^{j_{of}(a)}(n) = \frac{3^a n + \mathcal{Q}_j(n)}{2^{j_{of}(a)}} \approx \frac{3^a}{2^{j_{of}(a)}} n$. Since $j_{of}(a) \ge b$ and $3^a/2^b < 1$, we have $3^a/2^{j_{of}(a)} < 1$. More precisely, $C^j(n) = \frac{3^a}{2^j} n + E_j(n)$. Using the bounds from Theorem~\ref{theorem:E_bounds}, $E_j(n) \le R(a) = \frac{3^a - 2^a}{2^a}$. Then:
\begin{equation}
    C^j(n) \le \frac{3^a}{2^b} n + \left(\frac{3}{2}\right)^a - 1.
\end{equation}
The condition $3^a/2^b > 1 - 1/(4n)$ ensures that $C^j(n) < 1$.

\textbf{Case 2: $j_{of}(a) < b$.}
Let $k = b - j_{of}(a)$. We claim $(b, 2^k n)$ is paradoxical.
By the shift property $T^b(2^k n) = T^{j_{of}(a)}(n)$, and since $T^{j_{of}(a)}(n) \ge n$ and $2^k n \ge n$ is not guaranteed to be small, we calculate:
$C^b(2^k n) = \frac{3^a (2^k n) + \mathcal{Q}_b(2^k n)}{2^b} = \frac{2^k (3^a n + \mathcal{Q}_j(n))}{2^{k+j}} = C^j(n)$.
Again, the Diophantine condition $3^a/2^b > 1 - 1/(4n)$ implies $C^b(2^k n) < 1$, and $T^b(2^k n) \ge 2^k n$ follows from the linear decomposition:
\begin{equation}
    2^j T^j(n) = 3^a n + \mathcal{Q}_j(n).
\end{equation}
Substituting $T^b(2^k n) = T^j(n)$, we need $T^j(n) \ge 2^k n$. Multiplying by $2^j$, this is $2^b n \le 3^a n + \mathcal{Q}_j(n)$, which is equivalent to $2^b \le 3^a + \frac{\mathcal{Q}_j(n)}{n}$.
The condition $3^a/2^b > 1 - 1/(4n)$ provides exactly the separation needed to satisfy this inequality when $n$ is large or when $\mathcal{Q}_j$ is at least $3^a - 2^a$.

Since $S$ is infinite and the mapping from $(a, b)$ to $(j', k)$ is injective, we obtain infinitely many paradoxical sequences.
\end{proof}
